\documentclass[12pt]{article}
\usepackage[utf8]{inputenc}
\usepackage[T2A]{fontenc}
\usepackage[ukrainian]{babel}
\usepackage[a4paper,margin=2cm]{geometry}
\usepackage{array}
\usepackage{enumitem}

\pagestyle{empty}

\begin{document}
\fontsize{12}{14}\selectfont

\begin{flushright}
Додаток Б
\end{flushright}

\begin{center}
\underline{Національний університет «Києво-Могилянська академія»}
\end{center}

\vspace{0.5cm}

\noindent Факультет {інформатики\hspace{10.3cm}}
\vspace{0.3cm}

\noindent Кафедра {комп'ютерних наук\hspace{9.3cm}}
\vspace{0.3cm}

\noindent Освітній ступінь {бакалавр\hspace{9.6cm}}
\vspace{0.3cm}

\noindent Спеціальність {122 Комп'ютерні науки\hspace{7.5cm}}
\vspace{0.3cm}

\noindent Освітня/освітньо-наукова програма {Комп'ютерні науки\hspace{8.6cm}}
\vspace{0.3cm}

\begin{flushright}
\begin{tabular}{l}
\textbf{ЗАТВЕРДЖУЮ} \\[0.3cm]
Завідувач кафедри Нагірна Алла Миколаївна \\[0.5cm]
«\rule{0.8cm}{0.4pt}» \rule{3cm}{0.4pt} 20\rule{0.8cm}{0.4pt} року
\end{tabular}
\end{flushright}

\vspace{0.5cm}

\begin{center}
\textbf{З А В Д А Н Н Я}

\textbf{ДЛЯ КУРСОВОЇ РОБОТИ СТУДЕНТУ}

\vspace{0.3cm}

{Алексєєнко Анастасії Дмитрівні\hspace{2.5cm}}
\end{center}

\vspace{0.3cm}

\noindent 1. Тема роботи {Оптимізація обчислень генетичних алгоритмів
\noindent {через паралельну обробку за допомогою CUDA\hspace{12.8cm}}

\vspace{0.5cm}
\noindent керівник роботи {Медвідь Сергій Олександрович, старший викладач\hspace{1cm}}

\vspace{0.5cm}
\noindent затверджені наказом вищого навчального закладу від «\rule{0.6cm}{0.4pt}» \rule{2.5cm}{0.4pt} 20\rule{0.6cm}{0.4pt} року № \rule{1.5cm}{0.4pt}}

\vspace{0.2cm}

\noindent 2. Строк подання студентом роботи \underline{до 24 березня 2026 року\hspace{1.5cm}}

\vspace{0.2cm}

\noindent 3. План роботи

\vspace{0.3cm}

\small
\begin{enumerate}[label=\arabic*., leftmargin=*, itemsep=0.1cm]
\item {Introduction}
\begin{itemize}[noitemsep, topsep=0pt]
    \item General hypothesis on CUDA acceleration of genetic algorithm fitness evaluation.
    \item Research motivation.
\end{itemize}

\item {Theoretical Background}
\begin{itemize}[noitemsep, topsep=0pt]
    \item \textit{Literature Review:}
    \begin{itemize}[noitemsep]
        \item GPU-CPU architectures for GA parallelization.
        \item CUDA model suitability for GA-specific computations.
        \item Prior GPU-GA studies on fitness evaluation speedups.
        \item Implementation challenges and scalability limits.
    \end{itemize}
    \item \textit{Theoretical Framework:}
    \begin{itemize}[noitemsep]
        \item Parallel fitness evaluation kernel design principles.
        \item Compute-communication tradeoff modeling.
        \item Expected speedup bounds from Amdahl's law application.
    \end{itemize}
    \item \textit{Methodology:}
    \begin{itemize}[noitemsep]
        \item CUDA tools for Python.
        \item Benchmark configurations for typical optimization problems.
        \item Profiling tools and performance metrics definition.
    \end{itemize}
\end{itemize}

\item {Practical and Experimental Part}
\begin{itemize}[noitemsep, topsep=0pt]
    \item System architecture and components.
    \item Implementation of parallel computations using CUDA.
    \item GPU memory usage optimization.
\end{itemize}

\item {Results and Conclusions}
\begin{itemize}[noitemsep, topsep=0pt]
    \item Analysis of experimental results.
    \item CPU vs GPU implementation performance comparison.
    \item Conclusions and perspectives for further research.
\end{itemize}
\end{enumerate}

\vspace{0.5cm}

\end{document}